\documentclass[conference]{IEEEtran}
\IEEEoverridecommandlockouts


\usepackage{cite}
\usepackage{amsmath,amssymb,amsfonts,amsthm}
\usepackage{mathtools}
\usepackage{graphicx}
\usepackage{textcomp}
\usepackage{xcolor}
\usepackage{microtype}
\usepackage[largesmallcaps]{kpfonts}
\linespread{1.04}

\newcommand{\set}[1]{\left\{#1\right\}}
\newcommand{\pa}[1]{\left(#1\right)}
\newcommand{\bra}[1]{\left[#1\right]}
\newcommand{\abs}[1]{\left|#1\right|}
\newcommand{\norm}[1]{\left\|#1\right\|}
\newcommand{\der}[2]{\frac{\partial #1}{\partial #2}}
\newcommand{\pder}[2]{\pa{\frac{\partial #1}{\partial #2}}}

\renewcommand{\bigskip}{\vspace{2\baselineskip}}
\renewcommand{\medskip}{\vspace{1\baselineskip}}


\swapnumbers
\newtheoremstyle{plainsc}{8pt}{8pt}{\normalfont\itshape}{}{\normalfont\bfseries\scshape}{.}{.5em}{}
\newtheoremstyle{definitionsc}{8pt}{8pt}{\normalfont}{}{\normalfont\bfseries\scshape}{.}{.5em}{}
\newtheoremstyle{remarksc}{8pt}{8pt}{\normalfont}{}{\normalfont\bfseries\scshape}{.}{.5em}{}

\theoremstyle{definitionsc}
\newtheorem{definition}{Definition}[section]
\newtheorem{example}[definition]{Example}

\theoremstyle{plainsc}
\newtheorem{theorem}[definition]{Theorem}
\newtheorem{proposition}[definition]{Proposition}
\newtheorem{lemma}[definition]{Lemma}
\newtheorem{corollary}[definition]{Corollary}

\theoremstyle{remarksc}
\newtheorem{remark}[definition]{Remark}


\makeatletter
\def\footnoterule{\kern-3\p@
  \hrule \@width 2in \kern 2.6\p@}
\makeatother













\begin{document}

\title{%
  {\huge\textls[40]{\scshape\MakeUppercase{Biological Waste Treatment, \\Hydrology, and Indoor Air Quality}}}\\[0.6ex]
  {\footnotesize\textls[50]{\scshape\MakeLowercase{Course Project: Applied Stochastic Analysis (I2M411/MAT4511)}}}
}


\author{
  \IEEEauthorblockN{
    Sivasankaran R S, 
    Ashlin Nelson, 
    Aneeth Kumaar, 
    Amal Abdullah, 
    Sarang Sudhakar
  }
  \IEEEauthorblockA{
  	\{sivasankaran23, ashlin23, aneeth23, amal23, sarang23\}@iisertvm.ac.in\\
  	\\
    \textit{IISER Thiruvananthapuram, Kerala}
  }
}

\maketitle

\begin{abstract}
Stochastic differential equations (SDEs) provide a natural framework for modeling environmental and biological processes subject to intrinsic variability, turbulent dispersion, and random external perturbations. In this report, we formulate and investigate continuous stochastic systems describing biological waste degradation, hydrological streamflow/pollutant transport, and indoor air contaminant dynamics. We construct It\^o diffusion formulations, identify physical interpretations for drift and diffusion operators, and analyze numerical simulations performed via the Euler--Maruyama integration scheme.
\end{abstract}

\begin{IEEEkeywords}
Stochastic Differential Equations, It\^o Calculus, Biological Waste Degradation, Hydrological Transport, Indoor Air Quality.
\end{IEEEkeywords}











\section{Motivation}

\begin{itemize}

\item\textbf{The Engineering Problem}

Currently, to battle wastewater that major cities produce on a daily basis in huge amounts, we use bio-wastewater treatment plants, which use anaerobic digestion as its process to breakdown the filth into biodegradable matter, which is more safer than releasing the polluted water directly into natural water streams. The problems being neither the volume of the incoming water, nor the percentage of filth is constant, and unpredictable. Adding onto it, we will have sharp discharge peaks in the mornings and evenings, random flooding, and unaccounted industrial dumps of chemical waste every so often. We also have to account for delicate working conditions of the fellow bacterial slaves we have in custody.
		
		
	\begin{itemize}
		\item\textbf{Washout:} If the incoming flow is too erratic, the bacterial colony gets flushed out at a greater rate than their reproduction rate.
		
		\item\textbf{Toxic Shock:} An unaccounted spike in the toxic contaminants and pollutants that do pH alterations will cause the bio digestive system to fail and die faster.
		
	\end{itemize}

\break

\begin{figure}[htbp] 
    \centering
    \includegraphics[width=1.0\linewidth]{wastetanks_flowchart.png} 
    \caption{Water Treatment Plant Flowchart.}
    \label{fig:waste-flowchart}
\end{figure}

\item\textbf{Architecture}
\begin{itemize}
	\item\textbf{Tank 1}
	
	This tank is a buffer. Its purpose is to contain the surges, smooth out the flow, damping and mixing.
	
	\item\textbf{Tank 2}
	
	This is the fixed volume tank which contains the live bacteria. It receives the stream from tank 1 and does the bio breakdown.
\end{itemize}

\break

\item \textbf{The need for Stochastic Calculus}

In general Chemical Engineering, these systems are modelled using ODEs which were derived from the equation of conservation of Mass:

\begin{equation*}
	\frac{d(\text{Mass})}{dt} = \text{MassIn} - \text{MassOut} - \text{MassUsed}
\end{equation*}

\medskip

Bio-degradation relies on non-linear kinetics, like the Monod Equation:

\begin{equation*}
	\mu(S) = \mu_{\text{max}}\frac{S}{K_{S}+S}
\end{equation*}

\medskip

This system assumes smooth deterministic flow of the pollutant concentration. Since this function is concave ($\mu''(S)<0$), the predicted value will usually be "better" than the measured value, and will only provide a "Best Case" scenario, which is basically useless as it is very unlikely to happen.

\medskip

Additionally, deterministic models practically fails to answer questions such as:

\begin{itemize}
	\item How long the plant withstand an erratic inflow before it fails?
	\item What is the probability that the effluent pollution exceeds environmental safety limits during a 5-day cycle?
\end{itemize}


\end{itemize}

\bigskip










\section{Derivations and Formulations}


This widely used model invented by CJ Harris \cite{cjharris} , assumes the wastewater treatment plant as a 5-dimensional state vector: $$X_t=[X_t^1, X_t^2, X_t^3, X_t^4, X_t^5]^T$$

\begin{table}[htbp]
    \centering
    \renewcommand{\arraystretch}{1.5}
    \begin{tabular}{l p{4.5cm} l}
        \hline
        \textbf{\shortstack[l]{Variable}} & \textbf{Physical Meaning} & \textbf{Unit} \\
        \hline
        $X_t^1$ & Volume of fluid in Tank 1 & $\mathrm{m^3}$ \\
        \hline
        $X_t^2$ & The flow rate of entering fluid & $\mathrm{m^3/day}$ \\
        \hline
        $X_t^3$ & Pollutant concentration of the inflow & $\mathrm{mg/L}$ or $\mathrm{g/m^3}$ \\
        \hline
        $X_t^4$ & Pollutant concentration flowing from Tank 1 to Tank 2 & $\mathrm{mg/L}$ \\
        \hline
        $X_t^5$ & Pollutant concentration flowing out of Tank 1 & $\mathrm{mg/L}$ \\
        \hline
    \end{tabular}
\end{table}



\subsection{\textbf{Equalization Tank Volume $X^1$}}

The mass balance of the wastewater tank is modeled from equation \ref{eq:mass_balance}:

\begin{equation}
\label{eq:mass_balance}
\frac{dV}{dt} = Q_{\text{out}} - Q_{\text{in}}
\end{equation}

Here, we set the optimal outflow rate from the tank 1 to be such below:

\begin{equation}
	Q^* = \frac{\alpha_1+1}{\beta_1}
\end{equation}

where, $\beta _{1}$ is the feedback gain coefficient, and $\alpha _{1}$ is a scaling constant.

\medskip

And the incoming flow is $X_t^2$, and if it exceeds this rate, the total derivative flips to negative, which shows the accumulation of water in the tank.

Hence we have the equation,

\begin{equation}
\label{eq:Xt1}
	dX_t^1 = \left(\frac{\alpha_1 + 1}{\beta_1} - X_t^2\right)dt
\end{equation}


\medskip







\subsection{\textbf{Inflow Rate $X^2$}}

The behaviour of this inflow rate is determined by two major forces: \textbf{Drift} (a collective trend), and \textbf{Diffusion} (environmental fluctuations or random noise).

The Drift term is just the constant growth, with a negative feedback, given by:

\begin{equation}
	\text{Drift Term}=\left(\alpha _{1}+1-\beta _{1}\frac{X_{t}^{2}}{X_{t}^{1}}\right)dt
\end{equation}

\medskip

When the inflow rate $X_t^2$ becomes too high with respect to the tank volume $X_t^1$, the feedback component ($-\beta_1 {X_t^2}/{X_t^1}$)  becomes larger than the baseline growth, hence it pulls back the flow rate towards balance.

\medskip

We use the Cox-Ingersoll-Ross (CIR) or Feller square-root diffusion model \cite{wfeller} for random disturbances. This stochastic noise is from standard Brownian Motion $dW_t^1$, scaled with the state ratio, given by:

\begin{equation}
	\text{Diffusion Term}=\sqrt{2\gamma _{1}\frac{X_{t}^{2}}{X_{t}^{1}}}dW_{t}^{1}
\end{equation}

\medskip

The boundary conditions at $X_t^2 \to 0$ when inflow stops\footnote{While the boundary condition of strict positive drift is a physical anomaly (since we literally cannot have zero inflow), it physically represents an immense pressure build-up somewhere upstream, and computationally ensures the non-negativity of the state variable during simulation.} are:

\begin{itemize}
	\item The diffusion coefficient becomes zero, so no random noise in the system.
	\item The drift remains strictly positive, $(\alpha_1 + 1)dt > 0$.
\end{itemize}

So combining them both, we get:
\begin{equation}
\label{eq:Xt2}
	dX_{t}^{2}=\left(\alpha _{1}+1-\beta _{1}\frac{X_{t}^{2}}{X_{t}^{1}}\right)dt+\sqrt{2\gamma _{1}\frac{X_{t}^{2}}{X_{t}^{1}}}dW_{t}^{1}
\end{equation}







\medskip






\subsection{\textbf{Inflow Dirtiness $X^3$}}




















































%\break

%\begin{thebibliography}{00}
%
%
%\bibitem{cjharris} HARRIS, C. J. (1977) ‘Modelling, simulation and control of stochastic systems with applications in wastewater treatment’, International Journal of Systems Science, 8(4), pp. 393–411. doi: 10.1080/00207727708942050.
%\bibitem{wfeller} Feller, W. (1951). Two singular diffusion problems. Annals of Mathematics, 54(1), 173–182
%
%\end{thebibliography}

\end{document}
